\documentstyle[11pt,graphicx,epsf,multirow,longtable]{y-jbook}
\topmargin      0.3cm
\textwidth      40zw   % 14.815 cm
\oddsidemargin  0.6cm
\evensidemargin 0.6cm
\headsep        0.8cm  % ヘッダと本文の間隔
\parindent 0.4cm
\pagestyle{myheadings}

\begin{document}
\pagenumbering{roman}%     目次はローマ数字
\begin{center}
  \thispagestyle{empty}
  \includegraphics[clip,width=1.6cm,height=1.8cm]{img/share/cit001.eps}\\
  \Large 平成28年度\quad 卒業論文\\
  \vspace{1.6cm}
  \gt \Huge 水素水を用いた人工知能のイデアに関する研究\\
  A Study on Idea of Artificial Intelligence using Hydrogen water\\
  \vspace{1.0cm}
  \mc \large  （指導教員：大川茂樹教授）\\
  \vspace{3.0cm}
  \gt \Large 千葉工業大学\quad 工学部\\
  \smallskip
  未来ロボティクス学科\\
  \vspace{0.5cm}
  \LARGE No.1326○○○\\
  \vspace{0.5cm}
  \Huge 大川\quad 太郎\\
  \huge OKAWA \quad Taro
\end{center}
\clearpage
\thispagestyle{empty}
\vspace*{2.4cm}

{\LARGE 本論文は，千葉工業大学工学部未来ロボティクス学科の}

\vspace{0.5cm}

{\LARGE 卒業論文として相応しいものと認めます．}\\

\vspace{9cm}

\begin{flushright}
{\LARGE 指導教員：\underline{\hspace{90mm}}}
\end{flushright}

\vspace{0.8cm}

\begin{flushright}
{\LARGE 平成\hspace{1.5cm}年\hspace{1.5cm}月\hspace{1.5cm}日}
\end{flushright}

% 	表紙ページは別ファイルで作成
\setlength{\baselineskip}{1.8zw}
\setlength{\textheight}{33\baselineskip}% ここでの設定は表紙には効かない
\addtolength{\textheight}{\topskip}     % この場所で設定する必要がある
% 行数は 34 行  topskip は文字領域の上端から
% 1 行目のベースラインまでの距離

% 図のレイアウト
\setcounter{topnumber}{5}%    ページ上部の図表は 5 個まで
\def\topfraction{1.00}%       ページの上 1.00 まで図表で占めて可
\setcounter{bottomnumber}{5}% ページ下部の図表は 5 個まで
\def\bottomfraction{1.00}%    ページの下 1.00 まで図表で占めて可
\setcounter{totalnumber}{10}% ページあたりの図表は 10 個まで
\def\textfraction{0.04}%      ページうち本文が占める割合の下限
%        これを 0 にすると本文が 1 行だけのページが出来る
%        0.04 にすると 1 行だけのページは防げる
%        0.1 くらいが良いかも知れない
\def\floatpagefraction{0.7}%  図表だけのページは少なくともこれだけで図表が占める

\chapter*{概要}
ここに概要を記入．
\addcontentsline{toc}{chapter}{概要}

\chapter*{Summary}
ここに Summary を記入．
\addcontentsline{toc}{chapter}{Summary}

%%%%%%%%%%%%%%%%%%%%%%%%%%%%%%%%%%%%%%%%%%%%%%%%%%%%%%%%%%%%%%%%%%%%%
% \pagenumbering{roman}%     目次はローマ数字
\tableofcontents % 目次作成
\clearpage

\pagenumbering{arabic} % アラビア数字
%%%%%%%%%%%%%%%%%%%%%%%%%%%%%%%%%%%%%%%%%%%%%%%%%%%%%%%%%%%%%%%%%%%%%

%%% 論文の中身を記載．
%%% 章ごとに tex ファイルを作り，inputで読み込むと見やすいかも
\chapter{序論}
ここは序論．水素水とは，爆発物である\footnote{水素とは，元素番号1の宇宙空間に最も多く存在する元素であり...}．


\chapter{理論}
ここは理論．
参考文献\cite{hogehoge}

\chapter{結論}
ここは結論．そう，結論．


\chapter*{謝辞}\markboth{謝辞}{謝辞}
感謝の気持ちを綴る．ありがとう．ありがとう．そして，ありがとう．
\begin{flushright}
平成29年2月　著者記す\\
\end{flushright}
\addcontentsline{toc}{chapter}{謝辞}


\addcontentsline{toc}{chapter}{参考文献}
\begin{thebibliography}{n}
% 参考文献を記載
\bibitem{hogehoge}参考文献
\end{thebibliography}
% \chapter*{付録}
% \addcontentsline{toc}{chapter}{付録}

\end{document}
